% ============================================================================
% Fichier  : partie0/P00_quarante_cinq_questions.tex
% Titre    : Prologue -- Quarante-Cinq Questions et Aucune Réponse Fiable
% Auteur   : MINKA MI NGUIDJOÏ Thierry Emmanuel
% Version  : 1.0 -- novembre 2026
% ============================================================================
% Source du cas : illustration LizbiZ de l'article A0, « Vers l'auditologie »
% (HAL hal-05701007), section 4. Faits d'A0 à ne pas modifier : grille de 45
% éléments en 5 domaines, domaine D incohérent, H ~ 1,2 bit, décomposition
% 5 x 5, usage stratégique de NEK, régression de 35 % à 18 mois, H ~ 0,6 et
% conformité de 52 % à 84 % six mois plus tard.
% Faits communs avec ACSI (MSN-2S3-MSI05), à garder identiques : P00 (13/2/5,
% binaire 16/20), C07 (D.1 à D.5 : 4/0/0, 1/1/2, 4/0/0, 2/0/2, 2/1/1),
% C08 (PT-14 : 31 départs, 19 / 7 / 5), C09 (A 16/3/1, B 8/3/1, C 1/10/1,
% E 8/6/6 ; novembre D 18/1/1, D_KL 0,26).
% Corrigés des exercices : livret séparé (corriges/Corr_P00.tex).
% ============================================================================

\chapter[Prologue : Quarante-Cinq Questions]{Prologue : Quarante-Cinq Questions et Aucune Réponse Fiable}
\label{chap:prologue}
\chaptermark{Prologue}

\epigraph{%
    « Une grille qui interdit de répondre \emph{je ne sais pas}\\
    ne supprime pas l'ignorance~: elle la déguise en certitude. »%
}{Thierry Minka, auditeur des systèmes d'information}

\niveauChapitre{Introductif}{%
    Aucun prérequis d'audit. Il suffit de savoir calculer une proportion et
    un logarithme. Ce prologue se lit avant toute définition~: il part
    d'une mission qui a échoué.%
}

\begin{boxobjectifs}
\begin{itemize}
    \item Reconstituer, à partir de ses traces, la chronologie d'une mission
          d'audit qui a échoué (opération LIZBIZ).
    \item Distinguer une impression, une constatation et une opinion
          d'audit, et reconnaître chacune dans un rapport.
    \item Identifier les défauts de méthode qui privent un rapport de toute
          valeur probante.
    \item Expliquer pourquoi une grille binaire \OUI{}/\NON{} produit des
          réponses incohérentes sur un système partiellement observable, et
          estimer le biais qu'elle introduit.
    \item Identifier les trois états de connaissance d'une réponse et
          utiliser la modalité \NEK{} (connaissance insuffisante).
    \item Calculer, pour un ensemble de réponses, le taux de conformité~$\mu$,
          le taux d'incertitude~$U$ et l'entropie~$H$, et savoir pourquoi on ne
          conclut jamais sur une seule de ces mesures.
    \item Énoncer les cinq réflexes de l'auditeur et le contenu minimal d'un
          papier de travail.
    \item Situer le cours~: sa méthode, ses trois cas, ses six parties, ses
          quinze travaux pratiques et le projet individuel.
\end{itemize}
\end{boxobjectifs}

\bigskip

% ============================================================================
\section*{Avant de commencer : un rapport que personne n'a voulu signer}
% ============================================================================

Douala, mardi 14~avril~2026, 9~heures. Dans la salle du conseil de
LiZbiZ~SA\footnote{La société LiZbiZ, ses dirigeants, ses prestataires, les
dates et les montants de ce cas sont fictifs. Les mécanismes d'audit, les
normes et les textes cités sont exacts. Le cas reprend l'illustration de
l'article \emph{Vers l'auditologie}~\cite{minka2026auditologie}, dont les
données sont conservées.}, le comité de direction a devant lui un rapport
d'audit de vingt-huit pages. Il a été commandé deux mois plus tôt, dans
l'urgence. Personne, autour de la table, ne veut le signer.

LiZbiZ est une entreprise de logistique. Elle stocke, achemine et livre des
marchandises pour le compte de ses clients, des importateurs et des
distributeurs, dans tout le pays. Son activité repose entièrement sur une
application intégrée, que les agences utilisent pour enregistrer les
entrées en entrepôt, préparer les expéditions, facturer et payer les
fournisseurs.

\begin{tableaucours}{L{4.0cm}L{8.2cm}}{%
    \tableauHeader{Élément} &
    \tableauHeader{Description}
}
Activité & Logistique~: entreposage, transport, livraison, pour le compte
de clients \\
Implantation & Siège à Douala~; douze agences régionales \\
Informatique & Quatre centres de données (Douala, Yaoundé, Garoua,
Bafoussam)~; une application intégrée commune à toutes les agences \\
Fonctions de l'application & Stocks et entrepôts~; expéditions~;
facturation~; achats et paiements des fournisseurs \\
Volume (exercice 2025) & 18\,340 expéditions livrées~; 4\,812 paiements
fournisseurs \\
Direction & M.~ETEME, directeur général~; Mme~NKENG, contrôleuse interne~;
M.~KOUNOU, responsable informatique \\
Commissaire aux comptes & Cabinet NDONGO \& ASSOCIÉS \\
\end{tableaucours}
\vspace{-0.8em}
\begin{center}{\small\itshape Tableau~P.1. Fiche d'identité de LiZbiZ}\end{center}

La mission était née de deux alertes. L'inventaire annuel des entrepôts a
fait apparaître un écart de \textbf{186~millions de francs CFA} entre le
stock enregistré dans l'application et le stock compté physiquement. Le
commissaire aux comptes a assorti sa certification des comptes d'une
réserve sur les stocks~: il n'a pas pu s'assurer que le chiffre inscrit au
bilan était juste. Le directeur général a alors demandé un audit «~de tout
le système d'information, et vite~».

L'équipe d'audit interne a travaillé avec une grille de \textbf{45~questions}
réparties en \textbf{5~domaines}, à cocher OUI ou NON. Les réponses ont été
recueillies auprès du responsable informatique, des chefs d'agence et des
responsables des centres de données. Le rapport conclut que le contrôle
d'accès à l'application est «~globalement maîtrisé~».

Le commissaire aux comptes en a demandé les papiers de travail. Il n'y en a
pas. Le responsable informatique conteste deux conclusions. Trois chefs
d'agence affirment n'avoir jamais été interrogés sur les points qu'on leur
attribue. Le comité refuse de valider le rapport.

\medskip

Ce prologue reconstitue cet échec, en cherche les causes, puis suit la
mission qui l'a repris. Il en tire les réflexes qui structurent tout le
cours. Il ne suppose aucune connaissance préalable~: les notions sont
nommées ici, et définies avec précision dans les chapitres qui suivent.

\begin{boxalert}
\textbf{Violation typique.} Conclure «~globalement maîtrisé~» sur un domaine
où les répondants se contredisent, sans avoir noté qui a répondu quoi, ni
sur quelle preuve.

\smallskip
\textbf{Règle~:} une conclusion d'audit qui ne peut pas être reliée à une
preuve conservée n'est pas une conclusion~: c'est une opinion personnelle.
\end{boxalert}

\separateur

% ============================================================================
\section{Ce qui s'est passé}
\label{sec:pro-chronologie}
% ============================================================================

\subsection{Reconstituer avant de juger}

Un ingénieur qui analyse une panne ne commence pas par chercher un
coupable~: il reconstitue la séquence des événements à partir des journaux.
L'auditeur fait de même avec une mission qui a échoué. Il rassemble les
traces disponibles (courriels, versions successives du rapport, comptes
rendus de réunion, fichiers de travail) et les ordonne. Il ne retient que ce
qu'une trace atteste~; ce qu'on lui raconte sans trace reste une
déclaration, à vérifier.

Mme~NKENG, la contrôleuse interne, a établi la chronologie du
Tableau~P.2 après le refus du comité. Chaque ligne y est rattachée à une
pièce.

\begin{tableaucours}{L{2.4cm}L{6.6cm}L{3.2cm}}{%
    \tableauHeader{Date} &
    \tableauHeader{Événement} &
    \tableauHeader{Pièce}
}
12 février 2026 & Le directeur général demande un audit « de tout le système
d'information » par courriel, sans objectifs, ni périmètre, ni référentiel
& Courriel du DG \\
16 février & L'équipe reprend une grille de 45 questions utilisée dans une
autre entreprise, en 5 domaines (A à E) & Fichier de la grille, propriétés
du document \\
17 février -- 13 mars & Réponses recueillies par téléphone et par
courriel~; l'identité des répondants n'est pas notée question par question
& Grille remplie (une seule colonne de réponses) \\
16 mars & Le domaine D (contrôle d'accès) présente des réponses
contradictoires~; l'équipe retient la réponse majoritaire & Courriel
interne à l'équipe \\
30 mars & Rapport remis~: contrôle d'accès « globalement maîtrisé » &
Rapport, version 1 \\
7 avril & Le commissaire aux comptes demande les papiers de travail~: il
n'en existe pas & Courriel du cabinet \\
9 avril & Le responsable informatique conteste deux conclusions par écrit
& Note de M.~KOUNOU \\
14 avril & Le comité de direction refuse de valider le rapport &
Procès-verbal du comité \\
\end{tableaucours}
\vspace{-0.8em}
\begin{center}{\small\itshape Tableau~P.2. Chronologie de la première
mission (reconstitution de la contrôleuse interne)}\end{center}

Cette chronologie ne contient que des faits. Elle ne dit pas encore
pourquoi la mission a échoué~: c'est l'analyse, qui vient ensuite et qui
doit être séparée des faits. Cette discipline, séparer ce qui est établi de
ce que l'on en déduit, est la première que l'auditeur apprend.

\subsection{Impression, constatation, opinion}

Le rapport de vingt-huit pages mélange trois sortes d'énoncés, qu'il faut
toujours distinguer.

\begin{boxdef}
\termecle{Impression}\index{impression}~: jugement formé sans démarche ni
preuve, souvent à partir d'un entretien, d'une visite ou d'une rumeur. Elle
oriente l'investigation~; elle n'entre jamais dans un rapport.

\medskip
\termecle{Constatation}\index{constatation} (\textit{finding})~: fait
vérifiable, établi par l'auditeur et appuyé sur une preuve conservée. Un
tiers compétent qui refait la même démarche sur les mêmes pièces doit
aboutir au même résultat.

\medskip
\termecle{Opinion d'audit}\index{opinion d'audit}~: conclusion motivée que
l'auditeur tire d'un ensemble de constatations, au regard d'un critère
connu à l'avance, et dont il assume la responsabilité.
\end{boxdef}

Le test le plus simple pour les distinguer tient en deux questions. Un tiers
peut-il refaire la démarche et trouver la même chose~? Si non, c'est une
impression. Si oui, l'énoncé dit-il un fait ou un jugement porté sur un
ensemble de faits au regard d'un critère~? Le fait est une constatation, le
jugement une opinion.

\begin{boxscenario}{Opération LIZBIZ. Le rapport}
\textbf{Extraits du premier rapport} (version du 30~mars), annotés par
Mme~NKENG.

\medskip
\begin{tableaucours}{L{6.6cm}L{5.6cm}}{%
    \tableauHeader{Extrait} &
    \tableauHeader{Annotation}
}
« Les équipes informatiques sont motivées et compétentes. » & Impression~:
aucune démarche, aucune preuve \\
« 9 questions sur 10 du domaine A reçoivent une réponse positive. » &
Constatation en apparence~; mais on ignore qui a répondu, et sur quelle
pièce \\
« Le contrôle d'accès est globalement maîtrisé. » & Se présente comme une
opinion~; aucun critère, aucune constatation conservée \\
« Les sauvegardes sont réalisées quotidiennement. » & Déclaration du
responsable informatique, reprise sans vérification \\
\end{tableaucours}

\medskip
\textbf{Analyse} (confiance élevée). Aucun énoncé du rapport n'est une
constatation au sens strict~: aucune réponse n'est rattachée à un
répondant, à une date ou à une pièce. Le rapport n'est pas faux~; il est
\emph{invérifiable}, ce qui, pour un rapport d'audit, revient au même.
\end{boxscenario}

\subsection{Quatre défauts de méthode}

L'analyse de la mission fait apparaître quatre défauts. Les trois premiers
sont des fautes de méthode classiques~; le quatrième est plus profond, et
c'est lui qui fonde ce cours.

\begin{enumerate}
    \item \textbf{Pas de mandat.} Un courriel n'est pas une lettre de
          mission. Sans objectifs, l'équipe ne savait pas ce que l'audit
          devait démontrer~; sans périmètre, elle ne savait pas où
          s'arrêter~; sans référentiel, elle ne savait pas à quoi comparer
          (Chapitre~\ref{chap:planifier}).
    \item \textbf{Pas de critère.} La grille vient d'une autre entreprise.
          Ses questions ne disent pas ce qui serait « conforme » chez
          LiZbiZ~: « les droits sont-ils revus régulièrement~? » ne fixe ni la
          fréquence, ni le responsable, ni la trace attendue
          (Chapitre~\ref{chap:decomposition}).
    \item \textbf{Pas de preuve.} Les réponses ne sont rattachées ni à des
          répondants, ni à des documents, ni à des traces du système. Aucun
          tiers ne peut refaire la démarche (Chapitre~\ref{chap:preuve}).
    \item \textbf{Un instrument inadapté.} Même conduite avec soin, avec un
          mandat, des critères et des preuves, la grille aurait échoué sur le
          domaine~D, parce qu'elle oblige à répondre OUI ou NON quand la
          réponse honnête est «~je ne sais pas~»
          (section~\ref{sec:pro-grille}).
\end{enumerate}

Le Tableau~P.3 rattache chacun de ces défauts au principe professionnel
qu'il viole. Les normes citées seront présentées au
Chapitre~\ref{chap:normes}~; il suffit ici de constater que les fautes de la
première mission ne sont pas des maladresses, mais des manquements à des
règles écrites.

\begin{tableaucours}{L{3.2cm}L{4.6cm}L{4.4cm}}{%
    \tableauHeader{Défaut} &
    \tableauHeader{Principe violé} &
    \tableauHeader{Source}
}
Pas de mandat & Une mission a des objectifs, un périmètre et des critères
définis avant de commencer & ISO~19011~\cite{iso19011_2018}~;
Normes de l'IIA~\cite{iia2024gias} \\
Pas de critère & L'opinion se forme au regard de critères d'audit connus
à l'avance & ISO~19011 \\
Pas de preuve & Approche fondée sur la preuve~: une conclusion repose sur
des informations vérifiables & ISO~19011 (principe d'approche fondée sur
la preuve) \\
Instrument inadapté & L'ignorance doit pouvoir être déclarée, mesurée et
traitée & A0 et A1~\cite{minka2026auditologie,minka2026nek} \\
\end{tableaucours}
\vspace{-0.8em}
\begin{center}{\small\itshape Tableau~P.3. Les quatre défauts de la première
mission et les principes qu'ils violent}\end{center}

% ============================================================================
\section{Pourquoi la grille a échoué}
\label{sec:pro-grille}
% ============================================================================

\subsection{Des réponses qui se contredisent}

Le Tableau~P.4 reproduit trois questions du domaine~D (contrôle d'accès),
telles que trois interlocuteurs y ont répondu lors de la première mission.

\begin{tableaucours}{L{5.6cm}ccc}{%
    \tableauHeader{Question (domaine D)} &
    \tableauHeader{Resp. informatique} &
    \tableauHeader{Chef d'agence} &
    \tableauHeader{Centre de Garoua}
}
Les comptes des agents partis sont-ils désactivés~? & OUI & NON & OUI \\
Les droits d'accès sont-ils revus chaque trimestre~? & OUI & NON & NON \\
Les mots de passe des comptes d'administration sont-ils changés~? & OUI & OUI & NON \\
\end{tableaucours}
\vspace{-0.8em}
\begin{center}{\small\itshape Tableau~P.4. Trois questions du domaine D et
leurs réponses binaires (extrait)}\end{center}

Qui ment~? Probablement personne. Chacun répond depuis sa position~:
\begin{itemize}
    \item le responsable informatique répond pour la \emph{procédure} qu'il a
          rédigée~: elle prévoit la désactivation des comptes, il répond donc
          OUI~;
    \item le chef d'agence répond pour ce qu'il \emph{voit} dans son
          agence~: un ancien magasinier a encore pu se connecter le mois
          dernier, il répond donc NON~;
    \item le responsable du centre de Garoua répond pour son \emph{site}, et
          sur la foi de ce qu'on lui a dit.
\end{itemize}
Interrogés plus tard, plusieurs d'entre eux reconnaissent qu'ils \emph{ne
savaient pas} vraiment. La grille ne leur laissait que OUI ou NON~: ils ont
choisi la réponse qui leur paraissait la plus probable, ou la plus
confortable.

\subsection{Trois états, deux cases}

Une réponse d'audit peut traduire trois états de connaissance~: le
répondant \emph{sait} que le contrôle est en place~; il \emph{sait} qu'il ne
l'est pas~; ou il \emph{ne sait pas}. Une grille binaire n'offre que deux
cases. Le troisième état est alors forcé dans l'une des deux, au hasard ou
par calcul. C'est l'observation de départ de l'article qui fonde
l'instrument de ce cours~\cite{minka2026nek}.

\begin{boxdef}
\termecle{Réponse NEK}\index{NEK (réponse)} (\textit{Not Enough Knowledge},
connaissance insuffisante)~: troisième modalité de réponse, distincte de
\OUI{} et de \NON{}, qui déclare que l'information disponible ne permet pas
de conclure. Elle n'est pas un aveu d'échec~: c'est une information sur
l'état de connaissance de l'organisation.

\medskip
\termecle{Instrument trichotomique}\index{instrument trichotomique}~: grille
dont chaque question admet les trois réponses \OUI{}, \NON{} et \NEK{}.
\end{boxdef}

\begin{boxalert}
\textbf{Violation typique.} Confondre \NEK{} et «~sans objet~». Une question
est sans objet quand le contrôle ne s'applique pas (une agence sans
entrepôt n'a pas d'inventaire)~; elle reçoit \NEK{} quand le contrôle
s'applique mais qu'on ignore s'il fonctionne.

\smallskip
\textbf{Règle~:} \NEK{} décrit l'état de connaissance~; «~sans objet~»
décrit le périmètre. Une question sans objet est retirée des calculs~; une
réponse \NEK{} y entre, et pèse.
\end{boxalert}

Les réponses contradictoires du domaine~D ne sont donc pas du bruit à
éliminer en retenant la majorité. Elles signalent que l'organisation
elle-même ne sait pas si ses contrôles d'accès fonctionnent. C'est
précisément ce qu'un auditeur doit découvrir, et ce que la grille binaire
cachait.

\begin{boiteRemarque}{Retenir la réponse majoritaire, est-ce raisonnable~?}
Non. Une majorité de OUI peut venir de répondants qui décrivent la
procédure écrite sans l'avoir vu appliquer. Une réponse minoritaire peut
venir du seul répondant qui a regardé la réalité. Un audit ne vote pas~: il
cherche la preuve. Le vote masque de plus la question la plus utile~:
\emph{pourquoi} les réponses divergent-elles~? Ici, parce que la
procédure existe au siège et n'est pas appliquée en agence~: c'est une
constatation bien plus précieuse que n'importe quelle majorité.
\end{boiteRemarque}

\subsection{Combien coûte la case forcée~?}

On peut estimer l'effet de la grille binaire. Lors de la seconde mission
(section~\ref{sec:pro-reprise}), les quatre mêmes interlocuteurs ont
répondu aux cinq questions du domaine~D deux fois~: d'abord avec la grille
binaire de la première mission, puis avec la grille trichotomique.

\begin{tableaucours}{L{4.4cm}cccc}{%
    \tableauHeader{Grille} &
    \tableauHeader{\OUI{}} & \tableauHeader{\NON{}} & \tableauHeader{\NEK{}} &
    \tableauHeader{Total}
}
Binaire (OUI ou NON) & 16 & 4 & · & 20 \\
Trichotomique (OUI, NON ou NEK) & 13 & 2 & 5 & 20 \\
\end{tableaucours}
\vspace{-0.8em}
\begin{center}{\small\itshape Tableau~P.5. Domaine D, mêmes questions, mêmes
répondants, deux grilles}\end{center}

Avec la grille binaire, la part des \OUI{} vaut $\hat\mu^{B} = 16/20 =
0{,}80$. Avec la grille trichotomique, elle vaut $\mu^{T} = 13/20 = 0{,}65$.
L'article A1 appelle la différence \emph{indice de réduction du
biais}~\cite{minka2026nek}~:
\[
    \rho = \hat\mu^{B} - \mu^{T} = 0{,}80 - 0{,}65 = 0{,}15.
\]
La grille binaire surestimait la conformité du domaine de quinze points.

D'où viennent ces quinze points~? Les cinq répondants qui, avec la grille
trichotomique, déclarent ne pas savoir ont dû, avec la grille binaire,
choisir une case. Trois ont choisi OUI, deux ont choisi NON. Si l'on note
$U$ la part des réponses \NEK{} et $\pi$ la part des ignorants qui
choisissent OUI quand on les y oblige, et si les répondants informés
donnent la même réponse avec les deux grilles, on a
\[
    \hat\mu^{B} = \mu^{T} + \pi\, U
    \qquad\text{donc}\qquad
    \rho = \pi\, U = \frac{3}{5} \times 0{,}25 = 0{,}15.
\]
Cette petite égalité, posée pour l'exemple, dit l'essentiel~: le biais de
la grille binaire est d'autant plus fort qu'il y a d'ignorance ($U$ grand)
et que les ignorants penchent vers la réponse confortable ($\pi$ proche
de~1). Or, dans un audit, la réponse confortable est presque toujours OUI.

\begin{boxinfo}
\textbf{Pourquoi OUI est-il la réponse confortable~?} Répondre NON, c'est
signaler une défaillance dont on peut être tenu pour responsable. Répondre
OUI ne coûte rien sur le moment. Ce n'est pas de la malhonnêteté~: c'est
une réponse rationnelle à une situation où l'on est interrogé par quelqu'un
qui pourrait vous reprocher la réponse. L'auditologie en fait un axiome,
l'\emph{interaction stratégique} entre auditeur et audité, étudié aux
Chapitres~\ref{chap:fraude} et~\ref{chap:incertitude}.
\end{boxinfo}

% ============================================================================
\section{La mission reprise}
\label{sec:pro-reprise}
% ============================================================================

En mai~2026, le comité de direction confie une seconde mission au cabinet
AUDIX CONSEIL. Le chef de mission, M.~BEYALA, obtient cette fois une lettre
de mission signée~: objectifs, périmètre (les quatre centres de données et
les douze agences), référentiels (la procédure interne de gestion des
accès, complétée par l'ISO/IEC~27002), livrables et calendrier. Il conserve
le domaine~D comme cible prioritaire et applique la démarche que ce cours
enseigne. Les résultats décrits ci-dessous sont ceux de l'illustration de
l'article \emph{Vers l'auditologie}~\cite{minka2026auditologie}.

\subsection{Cinq décisions de méthode}

\begin{enumerate}
    \item \textbf{Ajouter la modalité \NEK{}.} Chaque question admet
          désormais trois réponses, et chaque \NEK{} est justifié.
    \item \textbf{Mesurer l'incertitude.} Pour le domaine~D, l'équipe
          calcule l'entropie des réponses~: environ \textbf{1,2~bit}, sur un
          maximum de $\log_2 3 \approx 1{,}585$~bit. C'est le signe d'une
          forte opacité~: l'organisation ne sait pas ce qu'elle fait en
          matière d'accès. Reste à savoir si cette ignorance est sincère ou
          calculée.
    \item \textbf{Décomposer le domaine.} Le domaine~D est découpé en
          $5 \times 5$ sous-éléments, de sorte qu'aucun niveau ne demande à
          l'auditeur de tenir plus de sept éléments à l'esprit à la fois.
    \item \textbf{Comparer les contextes de réponse.} Les réponses données en
          entretien formel, où l'enjeu est élevé, diffèrent nettement de
          celles de l'auto-évaluation, où l'enjeu est faible~: certains
          répondants utilisent \NEK{} de façon stratégique.
    \item \textbf{Regarder dans le temps.} À partir de trois cycles d'audit
          antérieurs, une estimation de l'évolution du domaine~D donne
          \textbf{35~\%} de probabilité de régression dans les 18~mois.
\end{enumerate}

Les décisions~1 à~3 relèvent de la méthode enseignée dans la Partie~II~;
les décisions~4 et~5 ouvrent sur des travaux de recherche présentés dans la
Partie~VI. Le prologue s'attache aux trois premières.

\subsection{Trois mesures pour lire un domaine}

Le Chapitre~\ref{chap:conclure} définira trois mesures qui accompagnent
tout le cours. Elles se calculent sur un ensemble de réponses \OUI{},
\NON{} et \NEK{}~; on note $p(\OUI)$, $p(\NON)$ et $p(\NEK)$ leurs parts.

\begin{boxdef}
\termecle{Taux de conformité}\index{taux de conformité}~: $\mu = p(\OUI)$.

\medskip
\termecle{Taux d'incertitude}\index{taux d'incertitude}~: $U = p(\NEK)$.

\medskip
\termecle{Entropie}\index{entropie}~: dispersion des réponses, en bits,
\[
    H = -\sum_{m \in \{\OUI, \NON, \NEK\}} p(m) \log_2 p(m),
    \qquad 0 \leq H \leq \log_2 3 \approx 1{,}585,
\]
avec la convention $0 \log_2 0 = 0$.
\end{boxdef}

L'entropie vient de la théorie de l'information de
Shannon~\cite{shannon1948}. Elle vaut zéro quand toutes les réponses sont
identiques, et atteint son maximum quand les trois réponses sont également
représentées. Elle se mesure en bits~: un bit est l'information apportée
par la réponse à une question à deux issues également probables.

\begin{boxlecture}{Domaine D, second passage}
Cinq questions du domaine~D sont posées à quatre interlocuteurs~: 20
réponses, dont 13~\OUI{}, 2~\NON{} et 5~\NEK{}.

\medskip
\textbf{Constatation.} $\mu = 13/20 = 0{,}65$~; part des \NON{}~$= 0{,}10$~;
$U = 5/20 = 0{,}25$. L'entropie vaut
\[
    H = -(0{,}65\log_2 0{,}65 + 0{,}10\log_2 0{,}10 + 0{,}25\log_2 0{,}25)
    \approx 0{,}404 + 0{,}332 + 0{,}500 \approx 1{,}24~\text{bit}.
\]

\textbf{Analyse} (confiance moyenne). Avec 65~\% de \OUI{}, une grille
binaire aurait conclu à un domaine plutôt conforme. Mais un quart des
réponses déclarent ne pas savoir, et l'entropie est élevée~: on ne peut pas
conclure. La conformité apparente est \emph{opaque}~: c'est ce que l'article
\emph{Vers l'auditologie} appelle une \emph{non-conformité inconnue}, qui
appelle d'abord des investigations et une clarification des processus,
avant tout correctif technique (Chapitre~\ref{chap:conclure}).

\medskip
\hfill\lecture{0{,}65}{0{,}25}{1{,}24}
\end{boxlecture}

Pour un ingénieur, ces calculs s'automatisent en quelques lignes. Le
programme suivant, en Python, est celui qu'utilisent les travaux pratiques~;
il sera repris et complété au Chapitre~\ref{chap:conclure}.

\begin{lstlisting}[language=Python]
from math import log2

def mesures(oui, non, nek):
    """Taux de conformite mu, taux d'incertitude U, entropie H (bits)."""
    n = oui + non + nek
    parts = [oui / n, non / n, nek / n]
    h = -sum(p * log2(p) for p in parts if p > 0)   # 0 log 0 = 0
    return oui / n, nek / n, h

mu, u, h = mesures(13, 2, 5)
print(f"mu = {mu:.2f}  U = {u:.2f}  H = {h:.2f} bit")
# mu = 0.65  U = 0.25  H = 1.24 bit
\end{lstlisting}

\begin{boxinfo}
\textbf{Une règle à retenir dès maintenant.} On ne conclut jamais sur une
seule mesure. Une entropie nulle peut signifier une certitude, mais aussi
une ignorance unanime~: si tous les répondants répondent \NEK{}, l'entropie
est nulle et pourtant personne ne sait rien ($U = 1$). Le
Chapitre~\ref{chap:conclure} en tirera la règle de décision du cours~:
$U$ d'abord, $\mu$ ensuite, $H$ pour contrôler.
\end{boxinfo}

\subsection{Localiser l'incertitude}

Une mesure globale dit qu'il y a un problème~; elle ne dit pas où. La
troisième décision de M.~BEYALA, décomposer le domaine, répond à cette
question. Le domaine~D est découpé en cinq sous-domaines, et une question
par sous-domaine est posée aux quatre interlocuteurs.

\begin{tableaucours}{cL{4.2cm}cccc}{%
    \tableauHeader{Code} &
    \tableauHeader{Sous-domaine} &
    \tableauHeader{\OUI{}/\NON{}/\NEK{}} &
    \tableauHeader{$\mu$} & \tableauHeader{$U$} & \tableauHeader{$H$}
}
D.1 & Politique et organisation & 4 / 0 / 0 & 1,00 & 0,00 & 0,00 \\
D.2 & Cycle de vie des comptes & 1 / 1 / 2 & 0,25 & 0,50 & 1,50 \\
D.3 & Authentification & 4 / 0 / 0 & 1,00 & 0,00 & 0,00 \\
D.4 & Comptes à privilèges & 2 / 0 / 2 & 0,50 & 0,50 & 1,00 \\
D.5 & Journalisation et revue & 2 / 1 / 1 & 0,50 & 0,25 & 1,50 \\
\textbf{D} & \textbf{Ensemble} & \textbf{13 / 2 / 5} & \textbf{0,65} &
\textbf{0,25} & \textbf{1,24} \\
\end{tableaucours}
\vspace{-0.8em}
\begin{center}{\small\itshape Tableau~P.6. Le domaine D décomposé en cinq
sous-domaines}\end{center}

Lu d'un bloc, le domaine paraissait uniformément opaque. Décomposé, il
montre que la politique (D.1) et l'authentification (D.3) sont connues et
conformes, et que toute l'incertitude se concentre dans trois
sous-domaines~: la vie des comptes (D.2), les comptes à privilèges (D.4) et
les journaux (D.5). C'est là que les tests doivent porter. La
décomposition n'a rien appris de nouveau sur l'organisation~; elle a dit à
l'auditeur \emph{où chercher}.

\subsection{Du NEK au NON~: tester le contrôle D.2.1}

Le Chapitre~\ref{chap:decomposition} enseignera le gabarit canonique de
décomposition à huit niveaux~: domaine, sous-domaine, processus, activité,
contrôle, test, preuve, trace. Chaque contrôle y porte une \emph{fiche
d'assertion}, qui dit ce que l'auditeur affirme et comment il le vérifiera.
Voici celle du contrôle le plus incertain du sous-domaine D.2.

\begin{boxobjet}{D.2.1 · Comptes des agents partis}
\textbf{Position dans le gabarit~:} domaine D (contrôle d'accès)
$\rightarrow$ sous-domaine D.2 (cycle de vie des comptes) $\rightarrow$
processus «~départ d'un agent~» $\rightarrow$ activité «~retrait des
accès~» $\rightarrow$ contrôle D.2.1.

\textbf{Objectif~:} s'assurer qu'aucun agent parti ne conserve d'accès à
l'application intégrée.

\textbf{Critère~:} le compte d'un agent parti est désactivé dans les
48~heures qui suivent la date de départ enregistrée par les ressources
humaines.

\textbf{Question~:} le compte de chaque agent parti au cours de la période
a-t-il été désactivé dans les 48~heures~? Réponse \OUI{}, \NON{} ou \NEK{}.

\textbf{Test~:} rapprocher la liste des départs de la période et la liste
des comptes actifs, puis comparer les dates.

\textbf{Preuve et trace~:} extraction des départs signée par les ressources
humaines~; extraction des comptes actifs, avec son empreinte~; journal des
désactivations de l'application.

\textbf{Risque~:} un agent parti utilise son compte pour consulter ou
modifier des données, par exemple pour masquer un écart de stock.

\textbf{Conséquence~:} pertes non détectées, réserve du commissaire aux
comptes reconduite, manquement à l'obligation de sécuriser les données
personnelles des clients.
\end{boxobjet}

On remarque trois choses. La question n'a qu'une seule condition~: on peut
y répondre sans hésiter. La preuve attendue est nommée \emph{avant} le
test~: l'auditeur sait ce qu'il doit obtenir. Et la réponse \NEK{} est
prévue~: si les ressources humaines ne peuvent pas fournir la liste des
départs, l'auditeur ne répond ni OUI ni NON~; il constate qu'on ne peut pas
savoir, et il le consigne.

Sur D.2, les répondants étaient partagés~: un OUI, un NON, deux NEK. Les
déclarations ne permettaient pas de trancher. Le test, lui, le permet.

\begin{boxscenario}{Opération LIZBIZ. Tester D.2.1}
\textbf{Choix.} Les départs sont peu nombreux (31 de janvier à avril
2026)~: AUDIX CONSEIL teste la population entière plutôt qu'un
échantillon, en rapprochant la liste des départs fournie par les
ressources humaines et le journal des désactivations de l'application.

\medskip
\textbf{Constatation.} 31~départs~; 19 comptes désactivés sous
48~heures~; 7 désactivés plus tard~; 5 encore actifs à la date du test.

\medskip
\textbf{Cotation.} \NON{}~: le critère n'est pas satisfait, et la preuve
est établie sur la population entière.

\medskip
\textbf{Analyse} (confiance élevée). Les deux \NEK{} des répondants
n'étaient pas une faiblesse de l'enquête~: ils indiquaient exactement où la
connaissance manquait. Le test a transformé une incertitude en
non-conformité \emph{connue}, qui appelle un correctif précis.
\end{boxscenario}

Le test est consigné dans un papier de travail. Le
Chapitre~\ref{chap:preuve} en donnera la forme complète~; en voici
l'essentiel.

\begin{boxprocedure}
\textbf{PT-14 · Test du contrôle D.2.1}

\smallskip
\begin{tableaucours}{L{3.4cm}L{8.6cm}}{%
    \tableauHeader{Rubrique} &
    \tableauHeader{Contenu}
}
Mission, domaine & LiZbiZ, seconde mission~; domaine D, contrôle D.2.1 \\
Critère & Procédure interne de gestion des accès~: désactivation sous
48~heures \\
Population & Départs du 1\textsuperscript{er}~janvier au 30~avril 2026~:
31 agents (extraction RH signée) \\
Méthode & Rapprochement exhaustif des départs et du journal des
désactivations \\
Pièces & Extraction RH~; extraction des comptes actifs et journal, avec
leurs empreintes SHA-256 \\
Résultat & 19 sous 48~h~; 7 en retard~; 5 encore actifs \\
Cote & \NON{} \\
Réalisé par, revu par & Auditeur de l'équipe~; M.~BEYALA, chef de mission \\
\end{tableaucours}
\end{boxprocedure}

\subsection{Deux non-conformités, deux remèdes}

Le rapport de la seconde mission distingue deux situations, que la grille
binaire confondait.

\begin{tableaucours}{L{3.6cm}L{4.0cm}L{4.6cm}}{%
    \tableauHeader{Situation} &
    \tableauHeader{Signature} &
    \tableauHeader{Remède}
}
Non-conformité connue & Entropie faible, non-conformité élevée & Correctifs
techniques, contrôles renforcés \\
Non-conformité inconnue & Entropie élevée & Formation, clarification des
processus, nouvelles investigations \\
\end{tableaucours}
\vspace{-0.8em}
\begin{center}{\small\itshape Tableau~P.7. Deux types de non-conformité
(d'après~\cite{minka2026auditologie})}\end{center}

Six mois plus tard, en novembre~2026, l'entropie du domaine~D est descendue
à environ \textbf{0,6~bit}~: les mêmes questions donnent 18~\OUI{},
1~\NON{} et 1~\NEK{}. L'organisation sait désormais ce qu'elle fait en
matière d'accès. Dans le domaine où l'entropie était faible, la conformité
est passée de \textbf{52~\%} à \textbf{84~\%}. Le
Chapitre~\ref{chap:suivi} montrera comment comparer rigoureusement deux
cycles d'audit.

% ============================================================================
\section{Les cinq réflexes de l'auditeur}
\label{sec:pro-reflexes}
% ============================================================================

L'échec de la première mission et la réussite de la seconde se résument en
cinq réflexes. Chacun fait l'objet d'au moins un chapitre du cours.

\begin{tableaucours}{cL{4.2cm}L{6.0cm}c}{%
    \tableauHeader{N°} &
    \tableauHeader{Réflexe} &
    \tableauHeader{Ce qu'il évite} &
    \tableauHeader{Chapitre}
}
1 & Pas d'audit sans mandat écrit & Une mission sans objectif ni périmètre,
dont on conteste ensuite les conclusions & \ref{chap:planifier} \\
2 & Le critère avant le constat & Juger « maîtrisé » sans savoir ce qui
serait conforme & \ref{chap:decomposition} \\
3 & La preuve avant la conclusion & Des affirmations qu'aucun tiers ne peut
vérifier & \ref{chap:preuve} \\
4 & Le droit de répondre \NEK{} & Des certitudes forcées, des contradictions
tranchées par un vote & \ref{chap:preuve} \\
5 & Le contradictoire avant le rapport & Un audité qui découvre les
conclusions et les conteste & \ref{chap:mission} \\
\end{tableaucours}
\vspace{-0.8em}
\begin{center}{\small\itshape Tableau~P.8. Les cinq réflexes de
l'auditeur}\end{center}

Le cinquième réflexe mérite un mot, car la première mission l'a ignoré.
Le \termecle{contradictoire}\index{contradictoire} consiste à présenter
chaque constatation à l'audité \emph{avant} le rapport, à recueillir sa
réponse et à la consigner. L'audité peut produire une pièce que l'auditeur
n'avait pas, ou expliquer une cause. Lors de la seconde mission, M.~KOUNOU
a ainsi indiqué que les départs ne lui étaient pas notifiés par les
ressources humaines~: la cause du défaut de D.2.1 n'était pas la négligence
de l'informatique, mais l'absence de circuit d'information. La
recommandation en a été changée.

\begin{boxprocedure}
\textbf{Le minimum du papier de travail.} Pour qu'une réponse d'audit puisse
fonder une constatation, le papier de travail qui la porte indique au
moins~:
\begin{itemize}
    \item la question posée et le contrôle auquel elle se rattache~;
    \item l'identité et la fonction du répondant, la date et le mode de
          recueil (entretien, courriel, observation)~;
    \item la réponse \OUI{}, \NON{} ou \NEK{}, et la preuve qui l'appuie
          (document, extraction, observation)~;
    \item pour un \NEK{}, la raison pour laquelle on ne peut pas savoir~;
    \item la référence de la preuve conservée, et son empreinte s'il s'agit
          d'un fichier~;
    \item le nom de l'auditeur qui a recueilli la réponse, et celui du
          réviseur.
\end{itemize}
\end{boxprocedure}

\subsection{Ce que ce prologue dit à chaque métier}

Le cours prépare à six métiers. Chacun lit l'affaire LiZbiZ à sa manière.

\begin{boxmetier}{Auditeur SI}
\textbf{Reprendre une mission contestée}
\begin{enumerate}
    \item Obtenir un mandat écrit~: objectifs, périmètre, référentiels,
          livrables, calendrier.
    \item Ne pas réutiliser les réponses de la première mission~: elles ne
          sont rattachées à aucune preuve.
    \item Décomposer d'abord le domaine contesté, puis rédiger une fiche
          d'assertion par contrôle.
    \item Accepter et documenter les réponses \NEK{}~: elles disent où
          chercher.
    \item Présenter les constatations à l'audité avant le rapport, et
          consigner ses réponses.
\end{enumerate}
\end{boxmetier}

\begin{boxmetier}{Contrôleur interne}
\textbf{Ce que Mme~NKENG retient}

La contrôleuse interne n'audite pas~: elle conçoit et fait vivre les
contrôles. L'affaire lui montre qu'un contrôle écrit (la procédure de
désactivation) n'est pas un contrôle appliqué, et qu'un contrôle sans trace
(aucun journal des revues de droits) ne peut pas être démontré. Elle
inscrit à son plan la notification automatique des départs et la
conservation des comptes rendus de revue (Chapitre~\ref{chap:controle-interne}).
\end{boxmetier}

\begin{boxmetier}{Auditeur de sécurité}
\textbf{Un test plutôt qu'une déclaration}

Le rapprochement des départs et des comptes actifs est un test technique,
que l'auditeur de sécurité sait conduire sur des extractions et des
journaux. Il sait aussi qu'une extraction n'est une preuve que si l'on peut
établir sa source et son intégrité~: d'où l'empreinte SHA-256 du PT-14
(Chapitres~\ref{chap:preuve} et~\ref{chap:audit-technique}).
\end{boxmetier}

\begin{boxmetier}{RSSI}
\textbf{Le commanditaire qui doit savoir}

Le responsable de la sécurité des systèmes d'information est celui qui
aurait dû savoir si les comptes des agents partis étaient désactivés. Un
\NEK{} sur D.2 est un signal pour lui~: son propre dispositif de pilotage
ne lui donne pas l'information. Il reçoit le rapport, en discute les
constatations, et porte le plan d'action (Chapitre~\ref{chap:audit-smsi}).
\end{boxmetier}

\begin{boxmetier}{Gestionnaire de risques}
\textbf{Du registre à la réalité}

Le registre des risques de LiZbiZ classait déjà l'usage du compte d'un
agent parti parmi les risques intolérables. Le gestionnaire de risques
apprend que ce risque n'était pas seulement possible~: cinq comptes
restaient ouverts. L'audit lui fournit la donnée qui manquait pour coter
le risque résiduel (Chapitres~\ref{chap:risque}
et~\ref{chap:gestion-risques}).
\end{boxmetier}

\begin{boxmetier}{Responsable de la conformité}
\textbf{Une obligation légale en jeu}

Un compte actif d'agent parti donne accès aux données personnelles des
clients. Le responsable de la conformité y voit un manquement possible à
l'obligation de sécuriser ces données, au sens de la loi camerounaise sur
la protection des données à caractère personnel, et une question pour le
prochain audit de conformité (Chapitres~\ref{chap:droit}
et~\ref{chap:audit-conformite}).
\end{boxmetier}

% ============================================================================
\section{Le cours et son fil rouge}
\label{sec:pro-cours}
% ============================================================================

\subsection{Une méthode, une discipline}

Ce cours enseigne une méthode de conduite de la mission d'audit des
systèmes d'information, présentée intégralement au fil des chapitres. Elle
est l'instanciation opérationnelle, pour le praticien, du programme de
recherche de l'\termecle{auditologie}\index{auditologie}~: une science
appliquée de la réduction de l'incertitude dans l'audit des systèmes
d'information~\cite{minka2026auditologie}. Son instrument de réponse
\OUI{}~/ \NON{}~/ \NEK{} et ses mesures sont ceux de
l'article~\cite{minka2026nek}.

\begin{boxinfo}
\textbf{Ce que dit la recherche.} L'article \emph{Vers l'auditologie} (A0)
pose six axiomes, dont l'\emph{observabilité partielle}~: l'auditeur ne
voit jamais tout le système, il en voit des traces. L'article sur
l'incertitude épistémique (A1) en tire l'instrument trichotomique, ses
mesures et un protocole de validation. Le cours s'appuie sur ces deux
textes déposés, et sur des travaux publiés de théorie de l'information, de
théorie des jeux et de psychologie cognitive~; la Partie~VI y revient
(Chapitres~\ref{chap:auditologie} et~\ref{chap:incertitude}).
\end{boxinfo}

\subsection{Trois cas}

Trois organisations fictives accompagnent le cours, du prologue aux
études de cas intégratrices.

\begin{tableaucours}{L{3.0cm}L{5.4cm}L{3.8cm}}{%
    \tableauHeader{Cas} &
    \tableauHeader{Situation} &
    \tableauHeader{Thèmes}
}
LiZbiZ (Douala) & Logistique, quatre centres de données, douze agences~;
écart de stock et réserve du commissaire aux comptes & Méthode, contrôle
interne, fraude, applications, sécurité \\
BANQUE SANAGA (Douala) & Banque commerciale après une panne de trois jours
de son système central & Contrôle interne bancaire (COBAC), changements,
SMSI, continuité, tiers \\
ANPP (Yaoundé) & Agence publique de paiement, auditée par la juridiction
des comptes & Audit public, conformité, projets \\
\end{tableaucours}
\vspace{-0.8em}
\begin{center}{\small\itshape Tableau~P.9. Les trois cas fil rouge}\end{center}

\subsection{Le parcours}

Le cours suit le cycle de la mission, puis l'applique aux grands domaines
de l'audit des systèmes d'information.

\begin{tableaucours}{L{3.6cm}cL{7.0cm}}{%
    \tableauHeader{Partie} &
    \tableauHeader{Chapitres} &
    \tableauHeader{Question traitée}
}
I. Fondements & 1 à 6 & Qu'est-ce qu'auditer~? Qu'est-ce qu'un contrôle,
un risque, une fraude~? \\
II. Méthode de la mission & 7 à 12 & Comment conduire une mission, de la
lettre de mission au suivi~? \\
III. Cadre normatif & 13 à 16 & Quelles normes, quels référentiels, quel
droit, quelle déontologie~? \\
IV. Pratique de l'audit & 17 à 26 & Comment auditer la gouvernance, les
applications, les projets, la sécurité, la continuité, les tiers, la
conformité~? \\
V. Défis avancés & 27 à 29 & Audit continu, intelligence artificielle,
passage à la preuve judiciaire \\
VI. Recherche & 30 à 32 & L'auditologie, ses mesures et ses questions
ouvertes \\
\end{tableaucours}
\vspace{-0.8em}
\begin{center}{\small\itshape Tableau~P.10. Les six parties du cours}\end{center}

\subsection{Cours, travaux pratiques, projet}

Le cours compte 90~heures~: 50~heures de cours magistral, 30~heures de
travaux pratiques et 10~heures de travaux dirigés.

\begin{itemize}
    \item \textbf{Les quinze travaux pratiques} (deux heures chacun) se
          déroulent dans le Cyber Range, sous le code \codeUE. Chacun met en
          œuvre un chapitre sur les pièces d'un des trois cas~: matrice des
          droits, analyse de risques, détection de fraude, plan d'audit,
          décomposition, échantillonnage, mesures, conformité, requêtes SQL,
          SMSI, configuration réseau, test d'intrusion, continuité, nuage,
          synthèse. Le mode d'emploi en donne la liste.
    \item \textbf{Le projet individuel} occupe les travaux dirigés. Chaque
          élève conduit seul une mission complète dans un environnement
          virtuel dédié, qui contient des défaillances à découvrir. Il
          franchit quatre jalons~: cadrage, décomposition, exécution,
          restitution. Le prologue en est le modèle~: la seconde mission de
          LiZbiZ est exactement ce que l'on attend de lui.
    \item \textbf{Les exercices} de fin de chapitre sont gradués
          (\niveaubase{} application, \niveaumoyen{} analyse,
          \niveauexpert{} expert). Leurs corrigés, expliqués et annotés, sont
          réunis dans un livret séparé~: on cherche d'abord, on vérifie
          ensuite.
\end{itemize}

\separateur

\begin{boxresume}
\textbf{Ce qu'il faut retenir de ce prologue~:}
\begin{itemize}
    \item \textbf{Un rapport sans papiers de travail n'est pas un rapport}~:
          une conclusion doit pouvoir être reliée à une preuve conservée.
    \item \textbf{Impression, constatation, opinion}~: seule la constatation
          est un fait vérifiable~; l'opinion se fonde sur des constatations
          et un critère connu à l'avance.
    \item \textbf{Quatre défauts} ont ruiné la première mission~: pas de
          mandat, pas de critère, pas de preuve, un instrument qui interdit
          de dire «~je ne sais pas~».
    \item \textbf{Une grille binaire force l'ignorance} dans OUI ou NON~:
          elle surestime la conformité ($\rho = \pi U$), et ses
          contradictions sont un signal, pas un bruit.
    \item \textbf{La réponse \NEK{}} déclare une connaissance insuffisante~;
          elle dit où chercher. Elle ne se confond pas avec «~sans objet~».
    \item \textbf{Trois mesures}, $\mu$, $U$ et $H$, lisent un ensemble de
          réponses~; on ne conclut jamais sur une seule.
    \item \textbf{Décomposer localise l'incertitude}~; \textbf{tester la
          réduit}~: le NEK de D.2 est devenu un NON prouvé.
    \item \textbf{Cinq réflexes}~: mandat écrit, critère avant constat,
          preuve avant conclusion, droit au \NEK{}, contradictoire avant le
          rapport.
\end{itemize}
\end{boxresume}

\begin{boxinfo}
\textbf{Points de contrôle.} Vous avez compris ce prologue si vous savez~:
\begin{itemize}
    \item classer un énoncé de rapport en impression, constatation ou
          opinion, et le justifier~;
    \item dire pourquoi la majorité des réponses ne tranche pas une
          question d'audit~;
    \item calculer $\mu$, $U$ et $H$ pour 13, 2 et 5 réponses, sans
          calculatrice pour les deux premiers~;
    \item expliquer pourquoi $H = 0$ ne suffit pas à conclure~;
    \item énumérer les rubriques minimales d'un papier de travail.
\end{itemize}
\end{boxinfo}

% ============================================================================
\section*{Exercices}
% ============================================================================

\begin{boxexo}[title={\ding{45}\quad Exercice~P.1. Impression, constatation ou opinion~? \niveaubase}]
Classez chacune des phrases suivantes et justifiez en une ligne.
\begin{enumerate}[(a)]
    \item « Le responsable informatique nous a paru très compétent. »
    \item « Sur les 31 agents partis entre janvier et avril 2026, 5 avaient
          encore un compte actif à la date du test (PT-14, extraction et
          empreinte jointes). »
    \item « La sécurité des accès est satisfaisante. »
    \item « Au regard de la procédure interne (désactivation sous
          48~heures), le contrôle de désactivation des comptes n'est pas
          efficace. »
    \item « Selon le chef d'agence de Garoua, les mots de passe ne sont
          jamais changés. »
\end{enumerate}
\end{boxexo}

\begin{boxexo}[title={\ding{45}\quad Exercice~P.2. Trois mesures \niveaubase}]
Calculez $\mu$, $U$ et $H$ pour chacun des ensembles de réponses suivants,
puis dites en une phrase ce que chaque ensemble signifie.
\begin{enumerate}[(a)]
    \item 20~\OUI{}, 0~\NON{}, 0~\NEK{}.
    \item 0~\OUI{}, 0~\NON{}, 20~\NEK{}.
    \item 7~\OUI{}, 7~\NON{}, 6~\NEK{}.
    \item 1~\OUI{}, 1~\NON{}, 2~\NEK{} (sous-domaine D.2).
\end{enumerate}
Lequel de ces ensembles montre qu'une entropie nulle ne suffit pas à
conclure~?
\end{boxexo}

\begin{boxexo}[title={\ding{45}\quad Exercice~P.3. Réécrire trois questions \niveaumoyen}]
Reprenez les trois questions du Tableau~P.4.
\begin{enumerate}[(a)]
    \item Pour chacune, indiquez le critère implicite et réécrivez-la de
          façon qu'elle ne contienne qu'une seule condition vérifiable.
    \item Indiquez, pour chacune, la preuve qui permettrait de répondre
          \OUI{} ou \NON{}.
    \item Décrivez une situation dans laquelle l'auditeur devrait répondre
          \NEK{}, et ce qu'il fait alors.
\end{enumerate}
\end{boxexo}

\begin{boxexo}[title={\ding{45}\quad Exercice~P.4. Le coût de la case forcée \niveaumoyen}]
Dans une autre entreprise, 40 réponses sur un domaine donnent, avec une
grille trichotomique, 18~\OUI{}, 10~\NON{} et 12~\NEK{}.
\begin{enumerate}[(a)]
    \item Calculez $\mu^{T}$ et $U$.
    \item Avec une grille binaire, on suppose que les répondants informés
          répondent de même et que 90~\% des ignorants choisissent OUI.
          Calculez $\hat\mu^{B}$ et $\rho$.
    \item Une direction exige au moins 70~\% de réponses positives pour
          déclarer un domaine « maîtrisé ». Que conclut-elle avec chaque
          grille~? Qu'en pensez-vous~?
    \item Pour quelle valeur de $\pi$ les deux grilles donnent-elles le même
          taux de conformité~? Pourquoi cette situation est-elle peu
          probable en audit~?
\end{enumerate}
\end{boxexo}

\begin{boxexo}[title={\ding{45}\quad Exercice~P.5. Les traces de la chronologie \niveaumoyen}]
Reprenez le Tableau~P.2.
\begin{enumerate}[(a)]
    \item Pour chaque événement, indiquez si la pièce citée suffit à
          l'établir, ou s'il faut une pièce complémentaire.
    \item Repérez, dans la chronologie, le moment où chacun des quatre
          défauts de méthode s'est produit.
    \item Quel événement aurait dû déclencher l'arrêt de la mission, et
          qu'aurait-il fallu faire alors~?
\end{enumerate}
\end{boxexo}

\begin{boxexo}[title={\ding{45}\quad Exercice~P.6. Rédiger un papier de travail \niveaumoyen}]
Le 22~mai 2026, vous interrogez M.~KOUNOU sur la revue trimestrielle des
comptes d'administration. Il répond~: « Elle est faite, mais je n'ai pas
les comptes rendus sous la main~; ils sont chez le responsable de la
sécurité, en congé jusqu'au 10~juin. »
\begin{enumerate}[(a)]
    \item Quelle réponse cotez-vous~? Justifiez.
    \item Rédigez le papier de travail correspondant, avec toutes les
          rubriques minimales.
    \item Que faites-vous le 10~juin~? Quelles sont les cotes possibles
          ensuite, et à quelles conditions~?
\end{enumerate}
\end{boxexo}

\begin{boxexo}[title={\ding{45}\quad Exercice~P.7. Programmer les mesures \niveaumoyen}]
La seconde mission a lu les cinq domaines de LiZbiZ~: A~(16, 3, 1),
B~(8, 3, 1), C~(1, 10, 1), D~(13, 2, 5) et E~(8, 6, 6), en nombre de
réponses \OUI{}, \NON{} et \NEK{}.
\begin{enumerate}[(a)]
    \item À l'aide de la fonction \texttt{mesures} du prologue, calculez
          $\mu$, $U$ et $H$ pour chaque domaine.
    \item Classez les domaines par $U$ décroissant, puis par $\mu$
          croissant. Les deux classements coïncident-ils~?
    \item Modifiez la fonction pour qu'elle accepte un quatrième argument,
          le nombre de questions sans objet, et qu'elle l'écarte des
          calculs. Pourquoi ne faut-il pas l'ajouter aux \NEK{}~?
\end{enumerate}
\end{boxexo}

\begin{boxexo}[title={\ding{45}\quad Exercice~P.8. Voter ou chercher~? \niveaumoyen}]
Sur la question « les sauvegardes du centre de Bafoussam sont-elles
testées chaque mois~? », cinq interlocuteurs répondent~: OUI (responsable
informatique), OUI (responsable du centre), OUI (technicien), NON
(prestataire de maintenance), NEK (contrôleuse interne).
\begin{enumerate}[(a)]
    \item Que conclurait une équipe qui retient la majorité~?
    \item Quelle est la réponse la plus informative, et pourquoi~?
    \item Proposez le test qui permet de trancher, la preuve attendue et la
          trace recherchée.
\end{enumerate}
\end{boxexo}

\begin{boxexo}[title={\ding{45}\quad Exercice~P.9. Concevoir la grille du domaine D \niveauexpert}]
En vous appuyant sur le Tableau~P.6, rédigez une grille trichotomique de
cinq questions, une par sous-domaine du domaine~D.
\begin{enumerate}[(a)]
    \item Chaque question est atomique et porte sur un critère que vous
          énoncez, avec sa source.
    \item Pour chaque question, indiquez la preuve attendue et ce qui
          justifierait une réponse \NEK{}.
    \item Expliquez pourquoi une seule question par sous-domaine ne suffit
          pas pour conclure sur ce sous-domaine, et ce qu'il faudrait faire
          de plus.
\end{enumerate}
\end{boxexo}

\begin{boxexo}[title={\ding{45}\quad Exercice~P.10. La note au directeur général \niveauexpert}]
M.~ETEME vous demande pourquoi il devrait payer une seconde mission.
\begin{enumerate}[(a)]
    \item Rédigez une note d'une page qui explique, sans jargon, pourquoi le
          premier rapport ne peut pas être validé.
    \item Expliquez pourquoi 65~\% de réponses \OUI{} sur le domaine~D ne
          permettent pas de conclure que le domaine est maîtrisé.
    \item Proposez le contenu de la lettre de mission de la seconde
          mission~: objectifs, périmètre, référentiels, livrables,
          calendrier.
    \item Indiquez ce que le commissaire aux comptes attendra de cette
          seconde mission pour lever sa réserve.
\end{enumerate}
\end{boxexo}

\bigskip

\begin{boxinfo}
\textbf{Pour aller plus loin.}
\begin{itemize}
    \item L'article \emph{Vers l'auditologie}~\cite{minka2026auditologie},
          section~4, présente l'illustration LiZbiZ dont ce prologue est
          tiré.
    \item L'article sur l'incertitude épistémique~\cite{minka2026nek}
          définit l'instrument trichotomique, ses mesures et l'indice
          $\rho$.
    \item La norme ISO~19011~\cite{iso19011_2018}, article~4, énonce les
          principes de l'audit, dont l'approche fondée sur la preuve.
    \item L'article fondateur de Shannon~\cite{shannon1948} définit
          l'entropie~; sa première partie se lit sans difficulté.
\end{itemize}
\end{boxinfo}

\begin{boxinfo}
\textbf{Transition.} Le prologue a montré une mission qui échoue et une
mission qui réussit. Avant d'apprendre la méthode, il faut définir ce
qu'est un audit, ce qu'il garantit, et ce qu'il ne garantit pas. Le
Chapitre~\ref{chap:audit-auditologie} pose ces fondements et présente
l'audit comme une réduction de l'incertitude.
\end{boxinfo}
